% U3 6~7절 — 로봇 접촉 제어 전망 + 결론
\section{로봇 접촉 제어를 향해}
\label{sec:u3-outlook}

지금까지의 여정을 한 줄로 줄이면 이렇다.
임피던스는 힘과 흐름 사이의 동역학적 관계이고, 그 관계는 전기에서도
기계에서도 같은 2차 구조 --- 저장 두 개와 소산 하나 --- 로 나타난다.
로봇 접촉 제어는 이 언어의 자연스러운 다음 장면이다.
로봇 손끝에 가상의 스프링과 댐퍼를 달아 원하는 힘-운동 관계를 설계하는
임피던스 제어에서, 강성 \(k_d\)와 감쇠비 \(\zeta\)는 본 총설에서 준비한
바로 그 눈금으로 읽힌다 \cite{hogan1985impedancepart1}.
예를 들어 접촉에서 허용할 힘과 밀림을 정하면 강성의 출발값이 나오고,
목표점을 멀리 잡은 채 강성을 켜는 순간 \(\frac{1}{2}k_d\delta^2\)의
에너지가 이미 장전된다는 계산은 5절의 스프링 에너지 감각 그대로이다.

이 다음 장면 --- 임피던스 제어의 유도, 어드미턴스 제어와의 선택 기준,
그리고 실제 사고 사례에 기반한 안전한 파라미터 설정 원리 --- 는 본 총설의
2부에서 다룬다.
본 1부와 2부의 모든 수치 예제는 공개 MuJoCo \cite{todorov2012mujoco} 교육
시뮬레이터에서 재현할 수 있다.
예를 들어 저장 에너지를 계산하지 않고 큰 강성을 명령했을 때의 발사 사고는
저장소의 \texttt{configs/lab01\_msd/f2\_launch\_high\_energy.yaml} 설정
하나로 안전하게 겪어볼 수 있다.
전체 유도를 담은 확장판 원고, 시뮬레이터 코드, 실행 가이드는 공개
저장소\footnote{\url{https://github.com/ycpiglet/manipulator-control-tutorial}}에서
제공한다.

\section{결론}
\label{sec:u3-conclusion}

본 총설 1부는 임피던스라는 언어를 유도 생략 없이, 하나의 표기로, 역사적
동기와 함께 정리했다.
장거리 전신선로의 신호 왜곡이라는 실제 문제에서 출발해, 라플라스 변환을
정의에서부터 유도하고, RLC 회로와 질량-스프링-댐퍼가 같은 2차 구조를
공유함을 계수 비교와 수치 검산으로 확인했다.
독자가 이 여정에서 얻어야 할 것은 공식 목록이 아니라, 고유진동수와
감쇠비라는 두 눈금으로 어떤 2차 시스템이든 읽어내는 감각이다.
2부는 이 감각을 로봇 손끝으로 옮겨, 접촉 제어 방식의 선택과 안전한
파라미터 설정을 다룬다 \cite{lynchpark2017modernrobotics}.
